% \iffalse meta-comment
%
% Copyright (C) 2017- by Yanshuo Chu <yanshuoc@gmail.com>
%
% This file may be distributed and/or modified under the
% conditions of the LaTeX Project Public License, either version 1.3a
% of this license or (at your option) any later version.
% The latest version of this license is in:
%
% http://www.latex-project.org/lppl.txt
%
% and version 1.3a or later is part of all distributions of LaTeX
% version 2004/10/01 or later.
%
% \fi
%
% \section{开题中期报告类的第三方依赖}
%
% 第三方宏包的集中加载，与 \file{src/flow/hit-thesis-deps.dtx} 是一对，
% 顺序上的道理见那边。
%
% \changes{v3.2a}{2026/08/12}{三段 \texttt{report-deps-a/b/c} 合成一段。
%   原先 \texttt{report-deps-c} 还在 \file{src/hit-report.dtx} 里被切成六段，
%   六段之间夹的全是注释，不是代码}
% 两个类的 deps 没有合并：它夹在两串单类文件中间，文件名必须是单类专属的，
% 否则 docstrip 的全局输入表排不出一致的先后。内容也只有七成相同——学位论文
% 要索引与子图，报告不要；报告有三处按校区学位分岔的条件加载。
%
% ^^A ======================================================================
% ^^A Module report-deps: 第三方宏包集中加载与基础配置（含 hyperlink/geometry 内联）（hit-report）
% ^^A ======================================================================
%    \begin{macrocode}
%<*report-deps>
%    \end{macrocode}
%
% 各种宏包
%    \begin{macrocode}

%% 各种宏包
%    \end{macrocode}
%
% \changes{v3.2a}{2026/08/12}{删掉这里的 \pkg{etoolbox}。
%   \file{src/pages/hit-report-bib.dtx} 也装了一次，两处留一处}
%
% \changes{v3.2a}{2026/08/12}{art-deps-a 的引擎检查改用 \cs{sys\_if\_engine\_xetex:F}。报错
%   文本里的空格在 \cs{ExplSyntaxOn} 下会被吞，所以写成 \texttt{\~{}}}
%    \begin{macrocode}

\hit@setup@math
%    \end{macrocode}
%
% 删除 \pkg{newtxtext}，由于最新一次更新 (\texttt{revision 62369}) 中与 \pkg{ctex} 宏包冲突，且无法解决，将正文字体替换为 TG TermesX 与 TG Heros.
% \changes{v3.0.18}{2022/05/02}{删除 \pkg{newtxtext}，由于最新一次更新 (\texttt{revision 62369}) 中与 \pkg{ctex} 宏包冲突，且无法解决，将正文字体替换为 TG Termes 与 TG Heros}
% \changes{v3.0.19}{2022/05/05}{将 TG Termes 字体替换为 TG TermesX 字体，与 \pkg{newtxmath} 宏包更适配，删除 \option{Scale} 选项}
% \changes{v3.2a}{2026/08/12}{英文字体三件套设置从 artcls 拆出至 art-fonts 模块}
%    \begin{macrocode}
%    \end{macrocode}
%    \begin{macrocode}
%    \end{macrocode}
%
%    \begin{macrocode}
\RequirePackage{graphicx}
\RequirePackage{pdfpages}
\includepdfset{fitpaper=true}
\hit@setup@list@packages
%    \end{macrocode}
%
% 列表间距不调，走 \hologo{LaTeX} 的默认值——本类的规范没有这条要求。取值变量留空即可，
% 机制与取值都在 \file{src/utils/hit-list.dtx}。
%    \begin{macrocode}
\hit@setup@list@spacing
\hit@setup@footnote@packages
%    \end{macrocode}
%
% \changes{v3.2a}{2026/08/12}{\pkg{xeCJKfntef} 留在这里，理由见
%   \file{src/flow/hit-thesis-deps.dtx} 同一条}
%
% \pkg{newtx} 更新至 v1.7 版本后与 \pkg{ctex} 宏包冲突, 需要手动引入 \option{nofontspec} 选项. 判断方法为是否存在 v1.7 更新后的新宏包 \pkg{newtx.sty}
% \changes{v3.0.17}{2022/02/23}{\pkg{newtx} 更新至 v1.7 版本后与 \pkg{ctex} 宏包冲突, 需要手动引入 \option{nofontspec} 选项. 判断方法为是否存在 v1.7 更新后的新宏包 \pkg{newtx.sty}}
%    \begin{macrocode}

\RequirePackage{xeCJKfntef}
\RequirePackage{longtable}
\hit@setup@table
\hit@setup@footnote
\hit@setup@publist
\RequirePackage{booktabs}
%    \end{macrocode}
%
% 重定义 \pkg{booktabs} 宏包的默认线宽，绘制三线表时不再需要显式设置，以实现更好的内容与格式分离
% \changes{v3.0.20}{2022/05/06}{重定义 \pkg{booktabs} 宏包的默认线宽，绘制三线表时不再需要显式设置，以实现更好的内容与格式分离}
%    \begin{macrocode}
\hit@setup@table@rules

%    \end{macrocode}
%
% \changes{v3.2a}{2026/08/12}{natbib/hyperref/url 配置从 art-deps 拆出至 art-hyperlink 模块（对齐 book-hyperlink）}
%    \begin{macrocode}
%    \end{macrocode}
% ^^A ======================================================================

%    \begin{macrocode}
%    \end{macrocode}
%
% \changes{v3.2a}{2026/08/12}{geometry 配置从 art-deps 拆出至 art-geometry 模块（对齐 book-geometry）}
%    \begin{macrocode}
%    \end{macrocode}
%    \begin{macrocode}

\bool_if:NT \g__hit_debug_bool
  {
    \RequirePackage { layout }
    \RequirePackage { layouts }
    \RequirePackage { lineno }
  }
\RequirePackage{tabularx}
\RequirePackage{xltabular}
\RequirePackage{flafter}
\RequirePackage{subcaption}%支持子图，必须在 ccaption 之前
\RequirePackage{ccaption} %支持双语标题
%    \end{macrocode}
%
% \changes{v3.2a}{2026/08/12}{深圳博士中期那三个包挪进 \file{src/pages/hit-report-cover.dtx}，
%   只有那一份封面调}
% \changes{v3.2a}{2026/08/12}{删掉这里的 \pkg{fancyhdr}，
%   \file{src/flow/hit-report-pagestyle.dtx} 早就装过了}
% \changes{v3.2a}{2026/08/12}{\pkg{changepage}、\pkg{multicol}、\pkg{placeins}、
%   \pkg{multirow} 从报告类删掉。前两个只有学位论文的索引页调，后两个类里不调；
%   报告要用写 \cs{hitusepackage}}
%
% \changes{v3.2a}{2026/08/12}{\cs{bicaption} 支持省略英文图表名的四参数写法}
% \pkg{ccaption} 的 \cs{bicaption} 是五个参数，第四个是第二条题注用的图表名，
% 用户得写成 \texttt{Fig.} 加一个负细空格才能把 \cs{fnum@figure} 插的空格抵掉。
% 这里让第四个参数可以省略，省略时取 \option{const} 族里的 \option{figure-prefix-en}
% 与 \option{table-prefix-en}。五个参数的写法照旧——那才是 \pkg{ccaption} 的标准形式，
% 换到别的模板不用改。
%    \begin{macrocode}
\hit@setup@caption

%    \end{macrocode}
% \changes{v3.0.21}{2022/05/11}{删除 \pkg{xltxtra} 宏包, 来实现规范样式的上标。\issue[125]}
% 设置深圳校区开题报告的正文行距与毕业论文模板一致
% \changes{v3.0.04}{2020/07/29}{设置深圳校区开题报告的正文行距与毕业论文模板一致}
% 此处修改与正文行距保持一致
% \changes{v3.0.05}{2020/09/13}{此处修改与正文行距保持一致}
% \changes{v3.1d}{2025/03/03}{将深圳校区硕士中期报告图注实现为“图 1-1 标题”的格式}
% \changes{v3.2a}{2026/08/12}{图注/表注编号条件块从 artcls 拆出至 art-floats 模块}
%    \begin{macrocode}
%    \end{macrocode}
%    \begin{macrocode}
%    \end{macrocode}
%
% 添加字号、行距等设置
%    \begin{macrocode}
%    \end{macrocode}
%
% 更改哈尔滨校区博士硕士开题节标题为小三号字体
% \changes{v3.0.08}{2020/09/28}{更改哈尔滨校区博士硕士开题节标题为小三号字体}
%
% 进行字号、行距等设置
%    \begin{macrocode}

%% 进行字号、行距等设置
%    \end{macrocode}
%
% \changes{v3.1d}{2025/03/03}{增加 \cs{parindent} 的值。\issue[201]}
% \changes{v3.2a}{2026/08/12}{章节标题格式从 artcls 拆出至 art-chapter 模块}
%    \begin{macrocode}
%    \end{macrocode}
%    \begin{macrocode}
\bool_lazy_all:nT
  {
    { \bool_if_p:N \g__hit_weihai_bool }
    { \bool_if_p:N \g__hit_master_bool }
    { \bool_if_p:N \g__hit_proposal_bool }
  }
  { \pagestyle { hit@weihai@headings@main } }
%    \end{macrocode}
% \changes{v3.1c}{2023/04/16}{为深圳校区硕士开题设置章节样式与页眉页脚样式}
% 为深圳校区硕士开题设置章节样式与页眉页脚样式。%
% \changes{v3.1c}{2023/04/16}{为深圳校区硕士中期设置章节样式与页眉页脚样式}
% \changes{v3.1c}{2023/04/16}{深圳校区硕士中期不换页}
% 为深圳校区硕士中期设置章节样式与页眉页脚样式。%
% \changes{v3.1d}{2025/03/03}{根据 \issue[228] 更新深圳硕士中期节标题字号和间距}
% \changes{v3.2a}{2026/08/12}{深圳硕士的 ctexset 与 fancypagestyle 被拆为两个独立条件块，分别放入 art-chapter 与 artcls，保持行为等价}
%    \begin{macrocode}
%    \end{macrocode}

%  \changes{v3.1d}{2025/03/03}{目录格式将在下文中设置}
%    \begin{macrocode}
%     \ifboolexpr{bool {hit@interim}}{}{\pagestyle{hit@shenzhen@headings@front}}
%    \end{macrocode}
%
% 中文封面
%    \begin{macrocode}
%    \end{macrocode}
%    \begin{macrocode}
%    \end{macrocode}
% 引用和参考文献格式
% \changes{v3.2a}{2026/08/12}{参考文献格式从 artcls 拆出至 art-bib 模块}
%    \begin{macrocode}
%</report-deps>
%    \end{macrocode}
%
% \Finale
